\documentclass[11pt]{article}
\usepackage[T1]{fontenc}
\usepackage[utf8]{inputenc}
\usepackage{amsmath,amsthm,mathtools}
\usepackage{newpxtext,newpxmath}
\usepackage[scaled=0.93]{sourcesanspro}
\usepackage[a4paper,margin=25mm,headheight=15pt,footskip=13mm]{geometry}
\usepackage{microtype}
\usepackage{booktabs,array,enumitem}
\usepackage{xcolor}
\usepackage[most]{tcolorbox}
\usepackage{fancyhdr,lastpage,titlesec}
\usepackage{hyperref}
\usepackage{listings}
\definecolor{Ink}{HTML}{16354A}
\definecolor{Teal}{HTML}{176B77}
\definecolor{Pale}{HTML}{EFF6F7}
\definecolor{Muted}{HTML}{566773}
\hypersetup{colorlinks=true,linkcolor=Teal,citecolor=Teal,urlcolor=Teal,
  pdftitle={A counterexample and a sharp spacing-three criterion for products of q-integers},
  pdfauthor={Research note prepared with ChatGPT},pdfsubject={Unimodality; q-integers; exact counterexample}}
\pagestyle{fancy}
\fancyhf{}
\fancyhead[L]{\sffamily\small\color{Muted}Q-INTEGER UNIMODALITY}
\fancyhead[R]{\sffamily\small\color{Muted}Counterexample and exact criteria}
\fancyfoot[C]{\sffamily\small\color{Muted}\thepage\ / \pageref*{LastPage}}
\renewcommand{\headrulewidth}{0.3pt}
\titleformat{\section}{\Large\sffamily\bfseries\color{Ink}}{\thesection}{0.7em}{}
\titleformat{\subsection}{\large\sffamily\bfseries\color{Ink}}{\thesubsection}{0.7em}{}
\titlespacing*{\section}{0pt}{2.2ex plus .7ex minus .2ex}{1.1ex}
\titlespacing*{\subsection}{0pt}{1.7ex plus .5ex minus .2ex}{.8ex}
\newtheoremstyle{research}{8pt}{8pt}{\itshape}{}{
  \sffamily\bfseries\color{Ink}}{.}{.5em}{}
\theoremstyle{research}
\newtheorem{theorem}{Theorem}[section]
\newtheorem{proposition}[theorem]{Proposition}
\newtheorem{lemma}[theorem]{Lemma}
\newtheorem{corollary}[theorem]{Corollary}
\theoremstyle{definition}
\newtheorem{definition}[theorem]{Definition}
\newtheorem{remark}[theorem]{Remark}
\newcommand{\qi}[1]{[#1]_q}
\newcommand{\qir}[2]{[#1]_{q^{#2}}}
\newcommand{\coef}[2]{[q^{#1}]\,#2}
\newcommand{\SU}{\textup{SU}}
\newcommand{\Z}{\mathbb Z}
\newcommand{\R}{\mathbb R}
\DeclareMathOperator{\supp}{supp}
\numberwithin{equation}{section}
\setlength{\parskip}{3.5pt}
\setlength{\parindent}{1em}
\setlist{itemsep=3pt,topsep=5pt}
\tcbset{resultbox/.style={enhanced,breakable,colback=Pale,colframe=Teal,
  boxrule=.6pt,arc=1.5mm,left=3mm,right=3mm,top=2mm,bottom=2mm,
  fonttitle=\sffamily\bfseries,coltitle=white}}
\lstset{basicstyle=\ttfamily\small,columns=fullflexible,keepspaces=true,
  breaklines=true,frame=single,rulecolor=\color{Muted},showstringspaces=false,
  xleftmargin=3pt,xrightmargin=3pt}
\emergencystretch=1.5em

\begin{document}
\thispagestyle{plain}
\begin{center}
  {\sffamily\small\bfseries\color{Teal}RESEARCH NOTE \quad / \quad 18 SEPTEMBER 2026}\par
  \vspace{7pt}
  {\sffamily\bfseries\color{Ink}\LARGE A counterexample and a sharp\par
  spacing-three criterion\par
  for products of \boldmath$q$\unboldmath-integers}\par
  \vspace{8pt}
  {\small Prepared with ChatGPT; self-contained proofs and exact computational checks}
\end{center}

\begin{abstract}
We attack Conjecture 5.4 of Connelly, Ito, Martinez, Shevchenko, and Yang,
\emph{Unimodality of $q$-Fibonomial coefficients for small cases},
arXiv:2605.12822v1. Its necessity assertion is false:
$(1+q)^6(1+q^3)$ is unimodal but violates the proposed condition.
We then prove the conjectured sufficient condition for every spacing $r\ge2$.
For spacing three we obtain a complete replacement: if no $a_i$ is divisible
by three, then
\[
 \left(\prod_i\qi{a_i}\right)\qir{b}{3}\text{ is unimodal}
 \quad\Longleftrightarrow\quad
 b\le 1+\sum_i\left\lfloor\frac{a_i}{3}\right\rfloor
       +2\left\lfloor\frac{\#\{i:a_i\equiv2\pmod3\}}{6}\right\rfloor.
\]
A failed inequality produces a specified first-half coefficient drop of
exactly one. We also establish a general necessary degree bound, settle the
spacing-two criterion, and prove the sharp two-shift threshold
$(1+q)^n(1+q^r)$ unimodal if and only if $n\ge r^2-3$.
The arguments are finite and algebraic; computations are supplementary.
The main $q$-Fibonomial unimodality conjecture is not resolved here.
\end{abstract}

\begin{tcolorbox}[resultbox,title=The counterexample at a glance]
\[
 (1+q)^6(1+q^3)
 =1+6q+15q^2+21q^3+21q^4+21q^5+21q^6+15q^7+6q^8+q^9.
\]
Here $r=3$, $b=2$, and $a_1=\cdots=a_6=2$.
No $a_i$ is divisible by $3$, and
$1+\sum_i\lfloor a_i/3\rfloor=1<2$.
Nevertheless, the displayed coefficient sequence is weakly unimodal.
\end{tcolorbox}

\noindent\textbf{Status.} The statements below are accompanied by proofs and
reproducible integer-arithmetic tests. They have not been independently
refereed or checked by a proof assistant. The counterexample refutes the
specified source statement; historical priority of the counterexample and
the stronger theorems has not been established by an exhaustive literature review.

\section{The selected conjecture and the result}
\label{sec:target}
For a positive integer $a$, put
\[
 \qi{a}=1+q+\cdots+q^{a-1}.
\]
Throughout, $q$ is a formal indeterminate, $r\ge2$, $k\ge1$, and
$a_1,\ldots,a_k,b$ are positive integers. We study
\begin{equation}\label{eq:product}
 P_{\boldsymbol a,b,r}(q)
   =\left(\prod_{i=1}^{k}\qi{a_i}\right)\qir{b}{r}.
\end{equation}
Its coefficients are nonnegative integers, its constant and leading
coefficients are one, and its degree is
\begin{equation}\label{eq:degree}
 D=S+r(b-1),\qquad S=\sum_{i=1}^k(a_i-1).
\end{equation}
Every factor is symmetric, so $P$ is symmetric.

The source of the attack is \cite[Conjecture 5.4, p.~14]{CIMSY}.
For brevity write its proposed condition as
\begin{equation}\label{eq:condition}
 \mathcal C_r(\boldsymbol a,b):\quad
 (\exists i\colon r\mid a_i)
 \quad\text{or}\quad
 b\le 1+\sum_i\left\lfloor\frac{a_i}{r}\right\rfloor.
\end{equation}
The conjecture asserts that \eqref{eq:condition} is sufficient, and also
necessary when $k\le3$ \emph{or} $r\le3$.
The distinction between ``or'' and ``and'' matters: the case
$k=6,r=3$ lies inside the asserted necessity range.
The source reports tests with $k\le5$, $r\le6$, and
$\max\{a_i,b\}\le15$ \cite[p.~14]{CIMSY}.

\begin{proposition}[Exact disproof]\label{prop:counterexample}
The necessity assertion in \cite[Conjecture 5.4]{CIMSY} is false.
\end{proposition}
\begin{proof}
Take $r=3$, $b=2$, and $a_i=2$ for $1\le i\le6$.
The binomial theorem gives the coefficient list
\[
 (1,6,15,21,21,21,21,15,6,1).
\]
It increases weakly to its central plateau and then decreases weakly.
On the other hand, $3\nmid2$ and
$b=2>1+6\lfloor2/3\rfloor=1$, so \eqref{eq:condition} fails.
Since $r=3$, the claimed necessity applies to this example.
\end{proof}

This is a counterexample to an auxiliary product criterion, not to the
main $q$-Fibonomial conjecture in \cite{CIMSY}.
Our stronger conclusions, proved below, are summarized in
Table~\ref{tab:results}. All appearances of ``unimodal'' allow equal
consecutive coefficients.

\begin{table}[htbp]
\centering
\small
\begin{tabular}{@{}p{.25\linewidth}p{.64\linewidth}@{}}
\toprule
Setting & Conclusion proved in this note \\
\midrule
Every $r\ge2$ & The condition $\mathcal C_r$ is sufficient. With no divisible
factor, unimodality also requires $b\le1+\lfloor S/r\rfloor$. \\
$r=2$ & The original condition $\mathcal C_2$ is necessary and sufficient. \\
$r=3$ & With no divisible factor, the exact upper bound on $b$ is
$1+\sum_i\lfloor a_i/3\rfloor+2\lfloor t/6\rfloor$,
where $t=\#\{i:a_i\equiv2\pmod3\}$. \\
Two shifted binomial rows & $(1+q)^n(1+q^r)$ is unimodal exactly when
$n\ge r^2-3$, for $r\ge2$, $n\ge0$. \\
\bottomrule
\end{tabular}
\caption{Mathematical results, rather than extrapolations from the finite tests.}
\label{tab:results}
\end{table}

\section{Elementary tools for symmetric unimodality}
\label{sec:tools}
A finite nonnegative sequence $(c_0,\ldots,c_D)$ is \emph{unimodal} if,
for some index $m$, it is weakly increasing up to $m$ and weakly decreasing
after $m$. It is \emph{symmetric} if $c_j=c_{D-j}$.
For a symmetric sequence, unimodality is equivalent to
\begin{equation}\label{eq:firsthalf}
 c_j-c_{j-1}\ge0\qquad
 0\le j\le\lfloor D/2\rfloor,
 \qquad c_{-1}=0.
\end{equation}
We abbreviate ``symmetric and unimodal'' by $\SU$.
All coefficient sequences are extended by zero outside their support.

A shifted sequence $q^sH(q)$ has its center shifted by $s$.
In a sum of shifted polynomials, symmetry is understood about the
\emph{specified common center}, not automatically half the largest
exponent of each summand. This convention is essential in the decomposition
used below.

\begin{lemma}[Centered intervals and closure]\label{lem:closure}
Products of $\SU$ polynomials with nonnegative coefficients are $\SU$.
Nonnegative sums of $\SU$ coefficient sequences with the same center are
also $\SU$.
\end{lemma}
\begin{proof}
If $F(q)=\sum_{j=0}^D f_jq^j$ is $\SU$, let
$\alpha_j=f_j-f_{j-1}$ for $0\le j\le\lfloor D/2\rfloor$, with
$f_{-1}=0$. Then $\alpha_j\ge0$ and
\begin{equation}\label{eq:intervals}
 F(q)=\sum_{j=0}^{\lfloor D/2\rfloor}
          \alpha_j q^j\qi{D-2j+1}.
\end{equation}
To verify the identity, the coefficient at $q^m$ on the right is
$\sum_{j=0}^{\min(m,D-m)}\alpha_j=f_{\min(m,D-m)}=f_m$.
Each summand is the indicator sequence of an interval centered at $D/2$.

For positive $u,v$, the coefficients of $\qi{u}\qi{v}$ at exponents
$0\le\ell\le u+v-2$ are
\[
 \min\{\ell+1,u,v,u+v-1-\ell\}.
\]
They form a symmetric triangular or trapezoidal sequence.
Apply \eqref{eq:intervals} to two factors of degrees $D$ and $E$.
Every product of two interval summands is $\SU$ with center $(D+E)/2$.
On the left of this center, each first difference is nonnegative, so the
same is true of their sum. Symmetry supplies the right half. This also
proves the assertion about sums with a common center.
\end{proof}

\begin{lemma}[A divisible factor]\label{lem:divisible}
If $r\mid a_i$ for some $i$, then $P_{\boldsymbol a,b,r}$ is unimodal for
every positive integer $b$.
\end{lemma}
\begin{proof}
Write $a_i=rA$. The identity
\[
 \qi{rA}=\qi r\qir A r
\]
shows that the relevant two factors are
\[
 \qi r\qir A r\qir b r.
\]
The polynomial $[A]_z[b]_z$ has an $\SU$ coefficient sequence in $z$.
Replacing $z$ by $q^r$ and multiplying by $\qi r$ repeats each coefficient
exactly $r$ times. Repetition preserves symmetry and unimodality.
The remaining factors are ordinary $q$-integers, so
Lemma~\ref{lem:closure} finishes the proof.
\end{proof}

\section{A centered identity proves the sufficient condition}
\label{sec:sufficiency}
The following identity removes one block of length $r$ from an ordinary
factor and one term from the spaced factor, without changing the center
of the residual sequence.

\begin{lemma}[Centered peeling]\label{lem:peeling}
For integers $a>r$ and $b\ge2$,
\begin{equation}\label{eq:peeling}
 \qi a\qir b r
 =\qi{a+r(b-1)}+q^r\qi{a-r}\qir{b-1}{r}.
\end{equation}
Both terms on the right, with zero padding, are symmetric about the
center of the polynomial on the left.
\end{lemma}
\begin{proof}
Multiply \eqref{eq:peeling} by $(1-q)(1-q^r)$. The identity becomes
\begin{align*}
 (1-q^a)(1-q^{rb})
 &= (1-q^{a+r(b-1)})(1-q^r)\\
 &\quad +q^r(1-q^{a-r})(1-q^{r(b-1)}),
\end{align*}
which follows on expansion. The left side of \eqref{eq:peeling} has
polynomial degree
\[
 d=a-1+r(b-1).
\]
The first term on the right has degree $d$. The unshifted residual
polynomial has degree $d-2r$, so its shift by $q^r$ has center
$r+(d-2r)/2=d/2$.
\end{proof}

\begin{theorem}[Sufficiency for every spacing]\label{thm:sufficient}
For every $r\ge2$ and every number $k$ of ordinary factors,
$\mathcal C_r(\boldsymbol a,b)$ implies unimodality of
$P_{\boldsymbol a,b,r}$.
\end{theorem}
\begin{proof}
A divisible factor is covered by Lemma~\ref{lem:divisible}. Assume there
is none, and write
\[
 Q=\sum_i\left\lfloor\frac{a_i}{r}\right\rfloor.
\]
We prove sufficiency of $b\le Q+1$ by induction on $b$.
For $b=1$, the spaced factor equals one and the product of ordinary
$q$-integers is $\SU$ by Lemma~\ref{lem:closure}.

Let $b\ge2$. Then $Q\ge1$, so at least one $a_i$ is greater than $r$:
equality is impossible because there is no divisible factor.
Apply \eqref{eq:peeling} to this factor and multiply by all the others.
The first summand is a product of ordinary $q$-integers, hence $\SU$.
In the residual product, $a_i$ is replaced by $a_i-r$ and $b$ by $b-1$.
No new divisible factor appears, the floor sum becomes $Q-1$, and
\[
 b-1\le Q=(Q-1)+1.
\]
The induction hypothesis makes the residual product $\SU$.
Lemma~\ref{lem:peeling} shows that the shifted residual and the first
summand have the same center, even after multiplication by the other
factors. Their sum is therefore $\SU$ by Lemma~\ref{lem:closure}.
\end{proof}

\begin{corollary}[Reduction to small residues]\label{cor:core}
Suppose no $a_i$ is divisible by $r$, and write
\[
 a_i=rA_i+s_i,\quad 1\le s_i\le r-1,\qquad Q=\sum_i A_i.
\]
If $b>Q+1$ and
\begin{equation}\label{eq:core}
 \left(\prod_i\qi{s_i}\right)\qir{b-Q}{r}
\end{equation}
is $\SU$, then $P_{\boldsymbol a,b,r}$ is $\SU$.
\end{corollary}
\begin{proof}
Apply \eqref{eq:peeling} a total of $Q$ times, subtracting $r$ from
some $a_i$ at each step. The spaced parameter never drops below two,
since its final value is $b-Q\ge2$. The final residual is
\eqref{eq:core}. Reversing the steps, each step adds an $\SU$ product of
ordinary $q$-integers at the same center as the shifted residual.
\end{proof}

Corollary~\ref{cor:core} is a sufficient reduction, not an assertion that
residue reduction preserves unimodality in both directions for every $r$.
That converse will be obtained for $r=3$ by a separate obstruction argument.

\section{Periodic coefficient differences and a necessary bound}
\label{sec:residues}
Let
\[
 F(q)=\prod_i\qi{a_i}=\sum_{j=0}^S f_jq^j,
 \qquad R_s=\sum_{\substack{0\le j\le S\\j\equiv s\ (r)}}f_j,
 \quad 0\le s<r.
\]
The $R_s$ are the total coefficient masses in the residue classes modulo
$r$. Put $\delta_s=R_s-R_{s-1}$, with subscripts interpreted modulo $r$.
As a formal power series, define
\begin{equation}\label{eq:u-general}
 U(q)=\frac{F(q)}{1-q^r}=\sum_{n\ge0}u_nq^n,
 \qquad u_n=\sum_{\ell\ge0}f_{n-r\ell}.
\end{equation}
As usual, coefficients outside their support are zero.

\begin{lemma}[The periodic tail]\label{lem:tail}
If $n\ge\max\{1,S-r+2\}$, then
$u_n-u_{n-1}=\delta_{n\bmod r}$.
If also $n<rb$, the same difference occurs in
$P=F\qir b r$:
\begin{equation}\label{eq:taildiff}
 \coef nP-\coef{n-1}P=\delta_{n\bmod r}.
\end{equation}
\end{lemma}
\begin{proof}
For the stated $n$, the next possible omitted indices in the residue
classes of $n$ and $n-1$ are $n+r>S$ and $n-1+r>S$.
Consequently $u_n=R_{n\bmod r}$ and $u_{n-1}=R_{(n-1)\bmod r}$.
Finally,
\[
 P(q)=(1-q^{rb})U(q),
\]
so its coefficients agree with those of $U$ below exponent $rb$.
\end{proof}

\begin{theorem}[A necessary degree bound for every spacing]\label{thm:degreebound}
Suppose no $a_i$ is divisible by $r$. If $P_{\boldsymbol a,b,r}$ is
unimodal, then
\begin{equation}\label{eq:necessarydegree}
 r(b-1)\le S,
 \qquad\text{equivalently}\qquad
 b\le1+\left\lfloor\frac Sr\right\rfloor.
\end{equation}
\end{theorem}
\begin{proof}
Let $\zeta$ be a primitive $r$th root of unity.
Every factor
\[
 \qi{a_i}\big|_{q=\zeta}=
       \frac{1-\zeta^{a_i}}{1-\zeta}
\]
is nonzero, so $F(\zeta)\ne0$. The $R_s$ cannot all be equal, since
otherwise $F(\zeta)=R_0(1+\zeta+\cdots+\zeta^{r-1})=0$.
The differences $\delta_s$ sum to zero and are not all zero, so at least
one is negative.

Choose an integer $n$ in the $r$-term interval
\[
 S-r+2\le n\le S+1
\]
whose residue has a negative difference. Such an $n$ is automatically
positive. Indeed, if the interval includes nonpositive integers then
$S\le r-2$; at $n=0$ the difference is $R_0-R_{r-1}=f_0>0$,
and for any negative $n$ in this interval both relevant residue masses
vanish.
Thus Lemma~\ref{lem:tail} applies.

Suppose for a contradiction that $L=r(b-1)>S$.
Then $n\le S+1<rb$. If $L\ge S+2$, we also have
$n\le S+1\le\lfloor(S+L)/2\rfloor$. Equation~\eqref{eq:taildiff}
gives a negative first-half difference, contradicting unimodality.

The remaining possibility is $L=S+1$. In that case
$S\equiv-1\pmod r$. Symmetry of $F$ implies $R_0=R_{r-1}$, hence
$\delta_0=0$. The selected negative difference cannot be at $n=S+1$,
whose residue is zero. Therefore $n\le S=\lfloor(S+L)/2\rfloor$,
and again \eqref{eq:taildiff} contradicts unimodality.
\end{proof}

Together, Theorems~\ref{thm:sufficient} and \ref{thm:degreebound} give a
useful sandwich. When there is no divisible factor,
\begin{equation}\label{eq:sandwich}
 b\le1+\sum_i\left\lfloor\frac{a_i}{r}\right\rfloor
 \ \Longrightarrow\ P\text{ unimodal}\
 \ \Longrightarrow\
 b\le1+\left\lfloor\frac{\sum_i(a_i-1)}r\right\rfloor.
\end{equation}
The residue classes decide what happens in the possible gap.

\section{Spacing two: the original criterion is exact}
\label{sec:r2}
\begin{theorem}[Complete spacing-two criterion]\label{thm:r2}
The polynomial $(\prod_i\qi{a_i})\qir b2$ is unimodal if and only if some
$a_i$ is even or
\[
 b\le1+\sum_i\left\lfloor\frac{a_i}{2}\right\rfloor.
\]
\end{theorem}
\begin{proof}
Sufficiency follows from Theorem~\ref{thm:sufficient}.
If all $a_i$ are odd, write $a_i=2A_i+1$ and $Q=\sum_iA_i$.
Then $S=2Q$, so Theorem~\ref{thm:degreebound} gives the same necessary
bound $b\le Q+1$.
\end{proof}

There is also an exact failure certificate. With all $a_i$ odd,
$F(-1)=1$, so the even residue mass exceeds the odd residue mass by one.
For $b\ge Q+2$, the index $n=S+1=2Q+1$ lies below $2b$ and no later than
the center $Q+b-1$. Lemma~\ref{lem:tail} gives
\[
 \coef{2Q+1}P-\coef{2Q}P=-1.
\]
This argument also covers factors $a_i=1$.

\section{Spacing three: a complete replacement theorem}
\label{sec:r3}
Assume initially that no $a_i$ is divisible by three. Define
\begin{equation}\label{eq:Qt}
 a_i=3A_i+s_i,\quad s_i\in\{1,2\},\qquad
 Q=\sum_iA_i,\qquad t=\#\{i:s_i=2\}.
\end{equation}
Then
\begin{equation}\label{eq:SQt}
 S=3Q+t.
\end{equation}

\begin{theorem}[Sharp spacing-three criterion]\label{thm:r3}
The polynomial $P_{\boldsymbol a,b,3}$ is unimodal if and only if
\begin{equation}\label{eq:r3criterion}
 (\exists i\colon3\mid a_i)
 \quad\text{or}\quad
 b\le1+Q+2\left\lfloor\frac t6\right\rfloor,
\end{equation}
where $Q,t$ are as in \eqref{eq:Qt} in the absence of a divisible factor.
If no $a_i$ is divisible by three and the inequality fails, write
$t=6h+v$, $0\le v\le5$, and $M=Q+2h$. At the explicit index
\begin{equation}\label{eq:dropindex}
 n_0=3M+\left\lfloor\frac{v+3}{2}\right\rfloor
\end{equation}
we have
\begin{equation}\label{eq:unitdrop}
 n_0\le\lfloor D/2\rfloor,\qquad
 \coef{n_0}P-\coef{n_0-1}P=-1.
\end{equation}
\end{theorem}

The proof has two independent parts. A binomial-prefix lemma proves
sufficiency. A six-case residue table proves necessity, including the
unit-drop certificate.

\subsection{The binomial-prefix lemma}
\begin{lemma}[Enough ordinary binary factors]\label{lem:binomialcomb}
For integers $B\ge1$ and $t\ge0$, the polynomial
\[
 (1+q)^t\qir B3
\]
is unimodal whenever $t\ge6\lfloor B/2\rfloor$.
\end{lemma}
\begin{proof}
For $B=1$ this is an ordinary binomial polynomial. Suppose $B\ge2$,
put $m=\lfloor B/2\rfloor$ and $K=6m$, and first take $t=K$.
Thus $B=2m$ or $B=2m+1$.
Let
\[
 f_j=\binom Kj,\qquad
 u_n=\sum_{\ell\ge0}f_{n-3\ell},\qquad d_n=u_n-u_{n-1},
\]
with $f_j=0$ outside $0\le j\le K$ and $u_n=d_n=0$ for $n<0$.
Let $R_0,R_1,R_2$ be the full residue masses of $(1+q)^K$.

Modulo $\Phi_3(q)=1+q+q^2$, one has $(1+q)^3\equiv-1$ and hence
$(1+q)^6\equiv1$. Since $6\mid K$,
\[
 R_0+R_1q+R_2q^2\equiv1\pmod{\Phi_3(q)}.
\]
Two polynomials of degree at most two congruent modulo $\Phi_3$ differ
by a constant multiple of $\Phi_3$. It follows that $R_0-R_1=1$ and
$R_1=R_2$. Since $R_0+R_1+R_2=2^K$,
\begin{equation}\label{eq:binomialresidues}
 R_0=\frac{2^K+2}{3},\qquad
 R_1=R_2=\frac{2^K-1}{3}.
\end{equation}
Consequently
\begin{equation}\label{eq:deltasK}
 (\delta_0,\delta_1,\delta_2)=(1,-1,0).
\end{equation}

Symmetry $f_j=f_{K-j}$ gives a useful reflection identity.
The terms of residue $n$ omitted from $u_n$ have indices
$n+3,n+6,\ldots$. Reflecting these indices about $K/2$ gives
\[
 u_n=R_{n\bmod3}-u_{K-n-3}\qquad(n\ge-1).
\]
Subtract the same identity at $n-1$. For $n\ge0$ this yields
\begin{equation}\label{eq:reflection}
 d_n=\delta_{n\bmod3}+d_{K-n-2}.
\end{equation}

For $0\le n\le K/2$,
\[
 d_n=\sum_{\ell\ge0}\bigl(f_{n-3\ell}-f_{n-1-3\ell}\bigr)\ge1.
\]
Every summand is nonnegative because the binomial coefficients increase
up to $K/2$: explicitly, $f_j/f_{j-1}=(K-j+1)/j>1$ for
$1\le j\le K/2$. The summand with $\ell=0$ is a positive integer.
For $K/2<n\le K-2$, the reflected index $K-n-2$ lies between zero and
$K/2$. Equations~\eqref{eq:deltasK}--\eqref{eq:reflection} therefore give
$d_n\ge-1+1=0$.
At the last two indices,
\[
 d_{K-1}=\delta_2=0,\qquad d_K=\delta_0=1.
\]
We have proved
\begin{equation}\label{eq:prefixmonotone}
 d_n\ge0\qquad(0\le n\le K).
\end{equation}

Now
\[
 (1+q)^K\qir B3=(1-q^{3B})\sum_{n\ge0}u_nq^n.
\]
If $B=2m$, its degree is $2K-3$, its left-middle index is $K-2$,
and $3B=K$. If $B=2m+1$, its degree is $2K$, its middle index is $K$,
and $3B=K+3$. In either case the coefficients through the center agree
with $u_n$, and their differences are nonnegative by
\eqref{eq:prefixmonotone}. Symmetry proves unimodality for $t=K$.
Finally, multiplication by $(1+q)^{t-K}$ preserves $\SU$, proving the
claim for every $t\ge K$.
\end{proof}

\subsection{Sufficiency of the corrected criterion}
\begin{proof}[Proof of Theorem~\ref{thm:r3}, sufficient direction]
A divisible factor is handled by Lemma~\ref{lem:divisible}.
Otherwise assume the right-hand inequality in \eqref{eq:r3criterion}.
For $b\le Q+1$, Theorem~\ref{thm:sufficient} applies.
For $b>Q+1$, use Corollary~\ref{cor:core} with $r=3$.
The residue product is
\[
 (1+q)^t\qir{b-Q}{3}.
\]
Set $B=b-Q$. The assumed inequality says
$B\le1+2\lfloor t/6\rfloor$, which implies
$t\ge6\lfloor B/2\rfloor$.
Lemma~\ref{lem:binomialcomb} proves that the residue product is $\SU$,
and Corollary~\ref{cor:core} lifts this back to the original product.
\end{proof}

\subsection{Necessity and the coefficient-drop certificate}
\begin{proof}[Proof of Theorem~\ref{thm:r3}, necessary direction]
Suppose no $a_i$ is divisible by three, and write
$t=6h+v$, $0\le v\le5$, and $M=Q+2h$.
Then $S=3M+v$ by \eqref{eq:SQt}.

Let $R_0,R_1,R_2$ be the residue masses of
$F=\prod_i\qi{a_i}$ modulo three. Complete blocks of three terms vanish
modulo $\Phi_3(q)=1+q+q^2$, so
\begin{equation}\label{eq:corecongruence}
 F(q)\equiv\prod_i\qi{s_i}=(1+q)^t
       \equiv(1+q)^v\pmod{\Phi_3(q)}.
\end{equation}
The degree-at-most-two residue polynomials of $F$ and $(1+q)^v$
are congruent modulo $\Phi_3$, so they differ by a constant multiple
of $\Phi_3$. Their corresponding residue differences are therefore equal.
The binomial theorem for $v=0,\ldots,5$ gives
Table~\ref{tab:residues}.

\begin{table}[htbp]
\centering
\begin{tabular}{@{}rrrrrrr@{}}
\toprule
$v$ & $\delta_0$ & $\delta_1$ & $\delta_2$
 & $e=\lfloor(v+3)/2\rfloor$ & $n_0\bmod3$ & $n_0-S$ \\
\midrule
0 &  1 & -1 &  0 & 1 & 1 &  1 \\
1 &  1 &  0 & -1 & 2 & 2 &  1 \\
2 &  0 &  1 & -1 & 2 & 2 &  0 \\
3 & -1 &  1 &  0 & 3 & 0 &  0 \\
4 & -1 &  0 &  1 & 3 & 0 & -1 \\
5 &  0 & -1 &  1 & 4 & 1 & -1 \\
\bottomrule
\end{tabular}
\caption{The residue class selected by $n_0=3M+e$ always has difference $-1$.}
\label{tab:residues}
\end{table}

If the bound in \eqref{eq:r3criterion} fails, then $b\ge M+2$.
Choose $n_0$ from \eqref{eq:dropindex}.
The table shows $n_0\ge S-1$ and $n_0\ge1$.
Moreover,
\[
 n_0\le3M+4<3M+6\le3b.
\]
Thus Lemma~\ref{lem:tail} applies at $n_0$, and the selected table entry
gives the exact difference $-1$.
Finally,
\begin{align*}
 \left\lfloor\frac{D}{2}\right\rfloor
 &=\left\lfloor\frac{3M+v+3(b-1)}{2}\right\rfloor\\
 &\ge\left\lfloor\frac{6M+v+3}{2}\right\rfloor
 =3M+\left\lfloor\frac{v+3}{2}\right\rfloor=n_0.
\end{align*}
The negative difference occurs in the first half of the symmetric
sequence, contradicting \eqref{eq:firsthalf}.
This proves both necessity and \eqref{eq:unitdrop}.
\end{proof}

\subsection{Consequences, minimality, and the missing sixth factor}
\begin{corollary}[A two-parameter infinite family]\label{cor:binomialexact}
For $t\ge0$ and $B\ge1$,
\begin{equation}\label{eq:binomialexact}
 (1+q)^t\qir B3\text{ is unimodal}
 \quad\Longleftrightarrow\quad
 t\ge6\left\lfloor\frac B2\right\rfloor.
\end{equation}
\end{corollary}
\begin{proof}
For $t\ge1$, take $t$ ordinary factors $a_i=2$ in
Theorem~\ref{thm:r3}, so $Q=0$. The condition becomes
$B\le1+2\lfloor t/6\rfloor$, equivalent to \eqref{eq:binomialexact}.
For $t=0$, the spaced polynomial is unimodal only for $B=1$;
for $B\ge2$ its first coefficients include $1,0$ and a later $1$.
\end{proof}

Every pair $B\ge2$, $t\ge6\lfloor B/2\rfloor$ supplies a counterexample
to the original necessity assertion: the source condition would require
$B\le1$. Thus the degree-nine example is not an isolated coincidence.

\begin{corollary}[Minimality in the spacing-three branch]\label{cor:minimal}
Among counterexamples to the original necessity assertion with $r=3$,
at least six nontrivial ordinary factors are required, and the total
polynomial degree is at least nine. Both bounds are attained by
$(1+q)^6(1+q^3)$.
\end{corollary}
\begin{proof}
A counterexample has no divisible factor and satisfies
\[
 Q+1<b\le Q+1+2\lfloor t/6\rfloor.
\]
Hence $t\ge6$, so there are at least six factors of residue two, each
nontrivial. Also $S=3Q+t\ge6$ and $b\ge2$, whence
$D=S+3(b-1)\ge9$. The displayed example realizes equality.
\end{proof}

For $k\le5$, one necessarily has $t\le5$, so the correction
$2\lfloor t/6\rfloor$ vanishes. The corrected theorem therefore agrees
exactly with the old spacing-three criterion throughout that range.
This explains the relevance of the source's reported cutoff of five
ordinary factors \cite[p.~14]{CIMSY}; it is not a numerical accident.
The minimality assertion here is restricted to the $r=3$ branch and is
not a claim about every possible spacing.

\begin{table}[htbp]
\centering
\small
\begin{tabular}{@{}lrrrrl@{}}
\toprule
Ordinary factors & $Q$ & $t$ & Old bound & Exact bound & Extra admissible $b$ \\
\midrule
$\qi2^6$  & 0 & 6  & 1 & 3 & $2,3$ \\
$\qi2^{12}$ & 0 & 12 & 1 & 5 & $2,3,4,5$ \\
$\qi5^6$ & 6 & 6 & 7 & 9 & $8,9$ \\
$\qi2^5\qi4$ & 1 & 5 & 2 & 2 & None \\
\bottomrule
\end{tabular}
\caption{The exact bound means the largest $b$ for which the product with $\qir b3$ is unimodal.}
\label{tab:examples}
\end{table}

The mechanism is arithmetic: modulo $\Phi_3$, a residue-two factor is
$1+q$, whose sixth power is one. Each complete group of six such factors
allows two additional terms in the spaced factor. In contrast, for
$r=2$ the only nonzero residue is one, so no analogous correction exists.

\section{An additional exact theorem: two shifted binomial rows}
\label{sec:twoshift}
The smallest counterexample suggests isolating the family with $b=2$.
For this family there is a closed threshold for every spacing.
The proof in this section is independent of the spacing-three residue
table; it provides a further check on the special value six.

\begin{theorem}[Sharp two-shift threshold]\label{thm:twoshift}
For integers $r\ge2$ and $n\ge0$,
\begin{equation}\label{eq:twoshift}
 (1+q)^n(1+q^r)\text{ is unimodal}
 \quad\Longleftrightarrow\quad n\ge r^2-3.
\end{equation}
\end{theorem}
\begin{proof}
Multiplication by $1+q$ preserves $\SU$ by Lemma~\ref{lem:closure}.
It is consequently enough to prove unimodality at $n=r^2-3$ and failure
at $n=r^2-4$. A unimodal polynomial at a smaller exponent would, on
multiplication by a suitable power of $1+q$, contradict the latter failure.
For $r=2$, the two boundary polynomials are respectively
$1+q+q^2+q^3$ and $1+q^2$, so the claim follows directly.
Assume henceforth that $r\ge3$.

For an exponent $n$, put $N=n+1$ and define
\begin{equation}\label{eq:H}
 H_n(j)=\binom nj-\binom n{j-1}
       =\frac{N-2j}{N}\binom Nj.
\end{equation}
Binomial coefficients outside their natural range are zero. Then
$H_n(N-j)=-H_n(j)$.
The coefficient difference at $j$ in $(1+q)^n(1+q^r)$ is
\begin{equation}\label{eq:twoshiftdiff}
 H_n(j)+H_n(j-r).
\end{equation}

\medskip\noindent\textbf{Unimodality at $n=r^2-3$.}
Here $N=r^2-2$ and the polynomial degree $n+r$ is odd.
For a first-half index $j\le N/2$, both terms of
\eqref{eq:twoshiftdiff} are nonnegative.
For the remaining first-half indices, define
\[
 s=N+r-2j,\qquad
 u=j-r=\frac{N-r-s}{2},\qquad
 v=N-j=\frac{N-r+s}{2}.
\]
Because $j>N/2$ and $j\le(n+r-1)/2$, the integer $s$ is even and
satisfies $2\le s<r$. Furthermore $0\le u<v<N/2$, so $H_n(u),H_n(v)>0$.
By reflection, \eqref{eq:twoshiftdiff} is $H_n(u)-H_n(v)$.
It remains to prove
\begin{equation}\label{eq:ratio}
 R_s:=\frac{H_n(u)}{H_n(v)}
 =\frac{r+s}{r-s}
   \frac{\binom N{(N-r-s)/2}}{\binom N{(N-r+s)/2}}
 \ge1.
\end{equation}

At $s=2$,
\[
 R_2=\frac{(r+2)(N-r)(N-r+2)}
           {(r-2)(N+r)(N+r+2)}=1,
\]
because the numerator minus the denominator is
\begin{equation}\label{eq:ratio-base}
 4N(N-r^2+2)=0.
\end{equation}
For an even $s$ with $s+2<r$, the adjacent ratio is
\begin{equation}\label{eq:ratio-step}
 \frac{R_{s+2}}{R_s}
 =\frac{(r+s+2)(r-s)(N-r-s)(N-r+s+2)}
        {(r+s)(r-s-2)(N+r+s+2)(N+r-s)}.
\end{equation}
All factors in the denominator are positive. Expanding the numerator
minus the denominator gives
\begin{equation}\label{eq:ratio-positive}
 4Nr(N-r^2+s^2+2s+2)=4Nr\,s(s+2)>0,
\end{equation}
where the equality uses $N=r^2-2$.
Starting with $R_2=1$, induction over the even values of $s$ proves
\eqref{eq:ratio}. All first-half differences are nonnegative, and
symmetry gives unimodality.

\medskip\noindent\textbf{Failure at $n=r^2-4$.}
Now $N=r^2-3$ and the polynomial degree is even.
At its central index
\[
 j=\frac{n+r}{2},\qquad u=j-r=\frac{N-r-1}{2},
\]
reflection turns \eqref{eq:twoshiftdiff} into $H_n(u)-H_n(u+1)$.
Using \eqref{eq:H} and
$\binom N{u+1}=\binom Nu(N-u)/(u+1)$, this difference is
\begin{align}
 H_n(u)-H_n(u+1)
 &=\frac{\binom Nu}{N}
   \left((r+1)-(r-1)\frac{N+r+1}{N-r+1}\right)\notag\\
 &=-\frac{4\binom Nu}{N(N-r+1)}<0.\label{eq:centralnegative}
\end{align}
The denominator is positive for $r\ge3$.
Thus unimodality fails at the previous exponent. The initial closure
argument proves the full threshold \eqref{eq:twoshift}.
\end{proof}

For $r=3$, Theorem~\ref{thm:twoshift} gives the threshold $n=6$,
in agreement with Corollary~\ref{cor:binomialexact} for $B=2$.
Historical priority of Theorem~\ref{thm:twoshift} is not asserted here;
its proof is included because it arose naturally from the counterexample
and strengthens the resulting family analysis.

\section{Exact computations and reproducibility}
\label{sec:verification}
The proofs above establish the infinite statements. The accompanying
\texttt{verify.py} separately checks finite cases using Python integers,
without floating-point arithmetic or external packages.
The run recorded in \texttt{verification\_results.json} passed every
check in Table~\ref{tab:tests}.

\begin{table}[htbp]
\centering
\small
\begin{tabular}{@{}p{.69\linewidth}r@{}}
\toprule
Check & Number of cases \\
\midrule
Independent schoolbook versus optimized multiplication & 1,000 \\
Centered peeling identity on an integer grid & 3,003 \\
Exhaustive spacing-three multiset products & 30,539 \\
Deterministic random spacing-three products, no divisible factor & 30,000 \\
Spacing-three products with a divisible factor & 2,000 \\
General-spacing sufficient condition & 10,000 \\
General-spacing necessary degree bound & 10,000 \\
Spacing-two necessary-and-sufficient criterion & 5,000 \\
Two-shift boundary polynomials, $2\le r\le60$ & 177 \\
Exact two-shift negative central-difference formula & 58 \\
Spacing-three binomial-comb threshold samples & 237 \\
\bottomrule
\end{tabular}
\caption{Finite checks actually executed. Some families overlap; the counts are not
numbers of distinct mathematical theorems or independent proofs.}
\label{tab:tests}
\end{table}

\subsection{What was checked}
The exhaustive spacing-three test considers all multisets of ordinary
parameters from $\{1,2,4,5,7,8\}$, with zero through eight factors,
and all $b$ from one through the corrected bound plus two.
The zero-factor cases check the harmless empty-product boundary.
In this finite collection, 938 products violate the old necessity
condition while satisfying the corrected theorem.
Every predicted nonunimodal case is also checked at the explicit index
\eqref{eq:dropindex}; the drop is exactly $-1$, not merely negative.

The separate random spacing-three experiment uses seed 1977, chooses
$k$ uniformly from $1,\ldots,16$, chooses each $a_i$ from the integers
$1,\ldots,69$ not divisible by three, and chooses $b$ uniformly from
one through the corrected bound plus eight. The seed is reset for this
experiment so it is independent of how many preliminary tests are run.

The small counterexample is expanded by three different routes:
sliding convolution, the binomial theorem, and direct enumeration of
all $2^7=128$ choices of six weights one and one weight three.
The generic unimodality checker tests the entire rise-then-fall pattern;
it does not assume symmetry. A second checker uses symmetry and
first-half differences.

A separate Wolfram Language evaluation, using the connected
Wolfram kernel reporting version \texttt{15.0.1}, confirmed the
counterexample coefficients, the peeling identity, the residue table,
the polynomial identities \eqref{eq:ratio-base} and
\eqref{eq:ratio-positive}, and the two-shift boundary tests for
$2\le r\le16$. The exact input and returned text are included as
\texttt{verify.wl} and \texttt{wolfram\_output.txt}.
This is an independent computational check, not proof-assistant
verification of the arguments.

\subsection{Reproduction commands and output}
Use Python 3.10 or later:
\begin{lstlisting}
python3 verify.py --full --out verification_results.json
\end{lstlisting}
The command writes a JSON run summary, the small counterexample
certificate, and a CSV table of spacing-three binomial thresholds.
It returns a nonzero exit code if any check fails.
For a shorter smoke test, omit \texttt{--full} and choose a different
output filename to preserve the recorded full run.

Build this report with the supplied \texttt{Makefile} or run
\begin{lstlisting}
latexmk -pdf -interaction=nonstopmode -halt-on-error report.tex
\end{lstlisting}
A TeX installation with the packages listed in the preamble is needed.
No custom font files or proprietary packages are distributed.

\subsection{Closed-form decision versus coefficient expansion}
Theorem~\ref{thm:r3} gives a decision procedure requiring $O(k)$
integer-arithmetic operations: test divisibility, compute $Q$ and $t$,
and compare $b$ with the bound. If the answer is negative,
\eqref{eq:dropindex} gives the witness index directly.
No expansion of the polynomial is needed to decide the question.
The actual bit complexity depends on the sizes of the input integers.

The positive proof uses induction and the centered identity.
Literally unrolling all $Q$ peeling steps can be long when the parameters
are encoded in binary; no claim of a uniformly short expanded
certificate is being made. The full finite verifier expands coefficients
because it is intended as an independent check, not as the fastest
implementation of the classification theorem.

\section{Scope and further targets}
\label{sec:scope}
The exact counterexample disposes of the stated necessity assertion.
The stronger results retain and prove its sufficient half, classify all
spacing-two and spacing-three products, and identify a precise residue
obstruction rather than a merely numerical discrepancy.
The proofs explain both the first counterexample and the complete family
that it begins.

The source's main $q$-Fibonomial question is different. This note neither
proves nor disproves that main conjecture, and it does not challenge the
source's established small-case theorems.
For $r\ge4$, a general classification of
$(\prod_i\qi{a_i})\qir b r$ is not obtained here. In particular, the
necessity question for two or three ordinary factors at those spacings
is not settled by the present arguments.

A concrete continuation is to combine the centered reduction with the
full vector of residue differences at larger $r$.
Lemma~\ref{lem:tail} gives exact obstruction indices whenever the
periodic tail enters the first half. Corollary~\ref{cor:core} gives
positive constructions from residue products.
For $r=3$, these two mechanisms meet exactly because
$(1+q)^6\equiv1\pmod{\Phi_3(q)}$.
For larger $r$, the residual factors $\qi{s_i}$ range over more than
$1$ and $1+q$, so the same one-parameter six-case reduction is no longer
available. Determining when the positive and negative tests meet is a
specific next task, not a conclusion asserted in this note.

\appendix
\section{A short standalone certificate}
The disproof itself does not depend on any of the stronger theorems.
The following elementary Python fragment checks it using only the
binomial theorem.
\begin{lstlisting}[language=Python]
from math import comb

c = [
    (comb(6, j) if 0 <= j <= 6 else 0)
    + (comb(6, j-3) if 0 <= j-3 <= 6 else 0)
    for j in range(10)
]
assert c == [1, 6, 15, 21, 21, 21, 21, 15, 6, 1]
assert c == c[::-1]
assert all(c[j-1] <= c[j] for j in range(1, 5))
a, b, r = [2]*6, 2, 3
assert not any(x % r == 0 for x in a)
assert b > 1 + sum(x//r for x in a)
assert len(a) <= 3 or r <= 3
\end{lstlisting}
The distributed full verifier uses explicit exceptions instead of
\texttt{assert}, so running Python with optimization does not disable its
checks. The fragment above is only a compact readable certificate.

\begin{thebibliography}{1}
\bibitem{CIMSY}
Brendan B. Connelly, Ezekiel Ito, Thomas C. Martinez,
Olha Shevchenko, and Kacey Yang.
\newblock \emph{Unimodality of $q$-Fibonomial coefficients for small cases}.
\newblock arXiv:2605.12822v1 [math.CO], 2026.
\newblock \href{https://arxiv.org/abs/2605.12822v1}{Version-pinned arXiv record};
\href{https://arxiv.org/html/2605.12822v1}{HTML};
\href{https://arxiv.org/pdf/2605.12822v1}{PDF}.
\newblock Conjecture 5.4 and the adjacent testing statement, p.~14.
\end{thebibliography}
\end{document}
